\chapter{Allergic Reaction}

\section*{Definition}
An allergic reaction is an immune response to a substance that the person's immune system recognizes as harmful. Reactions range from mild, such as itching or hives, to severe and life-threatening (anaphylaxis). Many immediate reactions are IgE-mediated, but not every allergic reaction uses the same mechanism, and anaphylaxis can occur without hives.

\section*{When to See a Doctor}
\begin{itemize}
\item Seek medical advice for a reaction that is new, recurrent, widespread, or difficult to control, or for swelling of the lips, eyelids, or face.
\item Call emergency services immediately for difficulty breathing, wheezing, throat tightness, stridor, hoarseness, tongue or throat swelling, fainting, confusion, pale or clammy skin, or rapidly worsening symptoms.
\item Also treat repeated vomiting, severe abdominal pain, or symptoms involving more than one body system after a likely allergen as possible anaphylaxis.
\item Do not wait for symptoms to become severe or for a rash to appear; anaphylaxis can occur without skin findings.
\end{itemize}

\section*{How It Develops}
Allergic reactions involve sensitization and re-exposure. During first exposure, allergen-specific IgE antibodies bind to mast cells and basophils. On re-exposure, the allergen cross-links IgE, triggering mast cell degranulation and release of histamine, leukotrienes, and prostaglandins. These mediators cause widening of blood vessels, increased vascular permeability, bronchoconstriction, and mucus secretion.

\section*{Causes}
\begin{itemize}
\item Foods: peanuts, tree nuts, shellfish, eggs, milk, wheat, soy
\item Medications: penicillins, cephalosporins, sulfonamides, NSAIDs, aspirin
\item Insect stings: bees, wasps, hornets, fire ants
\item Environmental: pollen, dust mites, mold, pet dander
\item Latex
\item Contact: poison ivy, nickel, fragrances
\end{itemize}

\section*{Related Symptoms}
\begin{itemize}
\item Itching, flushing, or hives (urticaria)
\item Swelling of the lips, eyelids, tongue, or other tissues (angioedema)
\item Sneezing, runny or blocked nose, cough, wheezing, or shortness of breath
\item Nausea, vomiting, abdominal pain, or diarrhea
\item Dizziness, fainting, fast heartbeat, or signs of shock
\item Severe symptoms affecting breathing, circulation, or multiple body systems (anaphylaxis)
\end{itemize}

\section*{Medical Evaluation}
\begin{itemize}
\item A detailed history of the timing, amount, route, suspected trigger, symptoms, treatment, and previous reactions. Photographs and packaging or ingredient information may help.
\item Examination during or after the reaction, including vital signs and assessment of the airway, breathing, circulation, and skin. Serum tryptase may be considered after suspected anaphylaxis in selected cases, but it does not replace treatment.
\item After recovery, an allergist may use skin-prick testing or blood tests for specific IgE directed at suspected allergens. A positive test indicates sensitization and is not, by itself, proof of clinical allergy.
\item A medically supervised oral food challenge may confirm or exclude food allergy when the history and testing remain uncertain. It must not be attempted at home.
\end{itemize}

\section*{Treatment Options}
\begin{itemize}
\item Intramuscular epinephrine in the mid-outer thigh is first-line for anaphylaxis. If symptoms do not improve or return, a second dose may be needed while emergency help is being obtained. Emergency monitoring is required because symptoms can recur.
\item A non-sedating antihistamine such as cetirizine may help itching or hives in a mild, skin-limited reaction; antihistamines do not treat airway obstruction, low blood pressure, or shock.
\item Corticosteroids are not a substitute for epinephrine and are not routinely relied upon for the initial treatment or prevention of biphasic anaphylaxis.
\item Bronchodilators may be added for persistent wheezing after epinephrine, but they do not replace epinephrine.
\item Allergen immunotherapy or other preventive treatment may be offered by an allergy specialist for selected environmental, venom, or medication-related allergies.
\end{itemize}

\section*{Self-Care and Home Management}
\begin{itemize}
\item For suspected anaphylaxis, use an epinephrine auto-injector immediately if available and call local emergency services; do not wait for an antihistamine to work.
\item Identify and avoid known allergens, but do not undertake unnecessarily broad elimination diets without medical advice.
\item Carry prescribed epinephrine auto-injectors if at risk for anaphylaxis, know how to use them, check their expiration dates, and follow a written emergency action plan.
\item Take an OTC antihistamine only for mild, skin-limited symptoms; it must not be used instead of epinephrine for breathing, throat, circulation, or severe gastrointestinal symptoms.
\item Wear medical alert identification
\item Inform healthcare providers of all allergies
\end{itemize}

\section*{References}
\begin{thebibliography}{99}
\bibitem{allergic_reaction:ref1} Cardona V, Ansotegui IJ, Ebisawa M, et al. World Allergy Organization anaphylaxis guidance 2020. \textit{World Allergy Organ J}. 2020;13:100472.
\bibitem{allergic_reaction:ref2} Sampson HA, Muñoz-Furlong A, et al. ``Second symposium on the definition and management of anaphylaxis.'' \textit{J Allergy Clin Immunol}. 2006;117(2):391-397.
\bibitem{allergic_reaction:ref3} Sicherer SH, Sampson HA. ``Food allergy: a review and update.'' \textit{J Allergy Clin Immunol}. 2018;141(1):41-58.
\bibitem{allergic_reaction:ref4} Adkinson NF, Bochner BS, et al. \textit{Middleton\'s Allergy: Principles and Practice}, 9th ed. Elsevier, 2020.
\bibitem{allergic_reaction:ref5} Shaker MS, Wallace DV, Golden DBK, et al. Anaphylaxis: a 2023 practice parameter update. \textit{Ann Allergy Asthma Immunol}. 2023;130(1):1--63.
\bibitem{allergic_reaction:ref6} Santos AF, Riggioni C, Agache I, et al. EAACI guidelines on the diagnosis of IgE-mediated food allergy. \textit{Allergy}. 2023;78(12):3057--3076.
\end{thebibliography}
